\documentclass[12pt,a4paper]{article}
\usepackage[utf8]{vietnam}
\usepackage{amsmath,amssymb}
\usepackage{enumerate}
\usepackage{tabularx}
\usepackage{tikz}
\usepackage{xcolor}
\usepackage[top=0.45cm, bottom=0.45cm, left=0.8cm, right=0.8cm]{geometry}
\DeclareUnicodeCharacter{00B2}{\ensuremath{^2}}
\DeclareUnicodeCharacter{00B3}{\ensuremath{^3}}

\setlength{\parskip}{0pt}
\setlength{\parindent}{0pt}
\setlength{\fboxsep}{3.5pt}
\setlength{\fboxrule}{0.5pt}

% Dòng kẻ chấm mờ cự ly chuẩn viết tay tiểu học (18pt)
\newcommand{\dongke}[1]{%
  \begingroup%
  \setlength{\parskip}{0pt}%
  \baselineskip=18pt%
  \vspace{0.4mm}%
  \foreach \i in {1,...,#1}{%
    \noindent{\color{gray!55}\dotfill}\par%
  }%
  \vspace{0.2mm}%
  \endgroup%
}

% Khung bao quanh bài tập tự luận
\newcommand{\khungBaiTap}[1]{%
  \vspace{0.8mm}%
  \noindent\fbox{%
    \parbox{\dimexpr\linewidth-2\fboxsep-2\fboxrule\relax}{%
      \vspace{0.3mm}%
      #1%
      \vspace{0.3mm}%
    }%
  }%
  \vspace{0.5mm}%
}

% Ô bài tập đặt tính
\newcommand{\oDatTinh}[3][4]{%
  \vspace{0.3mm}%
  \noindent\textbf{#2} #3\par%
  \dongke{#1}%
}

% Câu trắc nghiệm dàn 4 đáp án trên 1 hàng ngang
\newcommand{\cauTN}[6]{%
  \noindent\textbf{Câu #1.} #2\par\vspace{0.1mm}%
  \noindent%
  \makebox[0.25\linewidth][l]{\textbf{A.} #3}%
  \makebox[0.25\linewidth][l]{\textbf{B.} #4}%
  \makebox[0.25\linewidth][l]{\textbf{C.} #5}%
  \makebox[0.25\linewidth][l]{\textbf{D.} #6}\par\vspace{0.6mm}%
}

% Câu trắc nghiệm dàn 2 đáp án / hàng
\newcommand{\cauTNhaiDong}[6]{%
  \noindent\textbf{Câu #1.} #2\par\vspace{0.1mm}%
  \noindent%
  \makebox[0.5\linewidth][l]{\textbf{A.} #3}%
  \makebox[0.5\linewidth][l]{\textbf{B.} #4}\par\vspace{0.1mm}%
  \noindent%
  \makebox[0.5\linewidth][l]{\textbf{C.} #5}%
  \makebox[0.5\linewidth][l]{\textbf{D.} #6}\par\vspace{0.6mm}%
}

\begin{document}
\pagestyle{empty}
\fontsize{10.2pt}{13.2pt}\selectfont

% =============================================================================
% ========================= MÃ ĐỀ: 401 ========================================
% =============================================================================

% -------------------- TRANG 1: MẶT TRƯỚC --------------------
\noindent
\begin{tabularx}{\linewidth}{|X|c|}
\hline
\textbf{PHIẾU BÀI TẬP CUỐI TUẦN -- TOÁN LỚP 4 (TUẦN 11)} & \textbf{MÃ ĐỀ: 401} \\
\textit{Trung tâm Toán Thầy Lĩnh -- Năm học: 2026 - 2027} & \textit{Thời gian: 45 phút} \\
\hline
\end{tabularx}

\vspace{1mm}

\noindent
\begin{tabularx}{\linewidth}{|X|p{3.5cm}|}
\hline
\rule[-2mm]{0pt}{7.5mm}\textbf{Họ và tên:} \dotfill & \textbf{Lớp:} \makebox[2cm]{\dotfill} \\
\hline
\end{tabularx}

\vspace{1mm}

\noindent
\begin{tabularx}{\linewidth}{|p{3cm}|X|}
\hline
\centering\textbf{ĐIỂM SỐ} & \centering\textbf{LỜI NHẬN XÉT CỦA THẦY CÔ} \tabularnewline
\hline
\rule{0pt}{1.1cm} & \tabularnewline
\hline
\end{tabularx}

\vspace{1mm}

\noindent\textbf{BẢNG ĐIỀN ĐÁP ÁN TRẮC NGHIỆM (Mỗi câu đúng 0,25 điểm):}
\vspace{0.6mm}

\begin{tabularx}{\linewidth}{|c|*{12}{>{\centering\arraybackslash}X|}}
\hline
\textbf{Câu} & \textbf{1} & \textbf{2} & \textbf{3} & \textbf{4} & \textbf{5} & \textbf{6} & \textbf{7} & \textbf{8} & \textbf{9} & \textbf{10} & \textbf{11} & \textbf{12} \\
\hline
\textbf{Đáp án} & & & & & & & & & & & & \\
\hline
\end{tabularx}

\vspace{1.2mm}

\noindent\textbf{\large PHẦN I. TRẮC NGHIỆM KHÁCH QUAN (3,0 điểm)}

\cauTN{1}{Kết quả của phép tính $254\,618 + 138\,245$ là:}{$392\,863$}{$392\,853$}{$382\,863$}{$393\,863$}
\cauTN{2}{Biểu thức nào sau đây thể hiện tính chất giao hoán của phép cộng?}{$a + b = b + a$}{$(a+b)+c = a+(b+c)$}{$a \times b = b \times a$}{$a + 0 = a$}
\cauTN{3}{Giá trị của biểu thức $a + b + c$ với $a = 15\,000, b = 28\,000, c = 32\,000$ là:}{$65\,000$}{$70\,000$}{$75\,000$}{$80\,000$}
\cauTN{4}{Số thích hợp điền vào chỗ chấm $5\,\text{tấn}\,250\,\text{kg} = \dots\,\text{kg}$ là:}{$525$}{$5\,250$}{$50\,250$}{$52\,500$}
\cauTN{5}{Số thích hợp điền vào chỗ chấm $7\,\text{m}^2\,45\,\text{dm}^2 = \dots\,\text{dm}^2$ là:}{$745$}{$7\,045$}{$7\,450$}{$79$}
\cauTN{6}{Khoảng thời gian $3\,\text{phút}\,25\,\text{giây}$ bằng bao nhiêu giây?}{$185\,\text{giây}$}{$205\,\text{giây}$}{$215\,\text{giây}$}{$225\,\text{giây}$}
\cauTN{7}{Tìm số hạng chưa biết trong phép tính $x + 35\,420 = 89\,650$:}{$x = 54\,230$}{$x = 54\,130$}{$x = 125\,070$}{$x = 53\,230$}
\cauTN{8}{Điền dấu thích hợp: $45\,000 + 38\,000 \;\dots\; 50\,000 + 32\,000$}{$>$}{$<$}{$=$}{Không so sánh được}
\cauTN{9}{Số lớn nhất có $6$ chữ số khác nhau là:}{$999\,999$}{$987\,654$}{$987\,650$}{$986\,754$}
\cauTN{10}{Làm tròn số $476\,280$ đến hàng trăm nghìn ta được số:}{$400\,000$}{$470\,000$}{$480\,000$}{$500\,000$}
\cauTNhaiDong{11}{Một hình chữ nhật có chiều dài $45\,\text{m}$, chiều rộng $25\,\text{m}$. Chu vi hình chữ nhật đó là:}{$70\,\text{m}$}{$140\,\text{m}$}{$120\,\text{m}$}{$1\,125\,\text{m}^2$}
\cauTNhaiDong{12}{Một gia đình ngày thứ nhất ăn hết $800\,\text{g}$ gạo, ngày thứ hai ăn hết $1\,\text{kg}\,200\,\text{g}$ gạo. Cả hai ngày gia đình ăn hết số ki-lô-gam gạo là:}{$2\,\text{kg}$}{$20\,\text{kg}$}{$200\,\text{kg}$}{$1\,000\,\text{kg}$}

\vspace{0.4mm}
\noindent\textbf{\large PHẦN II. TỰ LUẬN (7,0 điểm)}

\noindent\textbf{Bài 1 (2,0 điểm). Đặt tính rồi tính:}
\vspace{0.4mm}

\noindent
\begin{tabularx}{\linewidth}{|X|X|X|X|}
\hline
\oDatTinh[4]{a)}{$345\,628 + 278\,195$} & 
\oDatTinh[4]{b)}{$482\,536 + 169\,487$} & 
\oDatTinh[4]{c)}{$528\,374 + 73\,648$} & 
\oDatTinh[4]{d)}{$614\,285 + 285\,715$} \tabularnewline
\hline
\end{tabularx}

\newpage
% -------------------- TRANG 2: MẶT SAU --------------------
\noindent
\begin{tabularx}{\linewidth}{|X|c|}
\hline
\textbf{PHẦN TỰ LUẬN (TIẾP THEO) -- PHIẾU BÀI TẬP TUẦN 11} & \textbf{MÃ ĐỀ: 401} \\
\hline
\end{tabularx}

\vspace{0.8mm}

% BÀI 2: TÌM X (4 Ý)
\khungBaiTap{%
  \textbf{Bài 2 (1,5 điểm). Tìm $x$:}\par\vspace{0.6mm}
  \noindent
  \begin{minipage}[t]{0.485\linewidth}
    \textbf{a)} $x - 45\,280 = 78\,350$\par
    \dongke{3}
    \vspace{0.4mm}
    \textbf{c)} $74\,500 - x = 28\,340$\par
    \dongke{3}
  \end{minipage}\hfill
  \begin{minipage}[t]{0.485\linewidth}
    \textbf{b)} $x + 36\,420 = 92\,580$\par
    \dongke{3}
    \vspace{0.4mm}
    \textbf{d)} $x : 8 = 14\,235$\par
    \dongke{3}
  \end{minipage}
}

\vspace{0.8mm}

% BÀI 3: ĐỔI ĐƠN VỊ (4 Ý - 1 CỘT TỪ TRÊN XUỐNG)
\khungBaiTap{%
  \textbf{Bài 3 (1,0 điểm). Viết số thích hợp vào chỗ chấm:}\par\vspace{0.8mm}
  \noindent
  a) $6\,\text{tấn}\,350\,\text{kg} = \dotfill\,\text{kg}$\par\vspace{1.4mm}
  \noindent
  b) $8\,\text{m}^2\,24\,\text{dm}^2 = \dotfill\,\text{dm}^2$\par\vspace{1.4mm}
  \noindent
  c) $5\,\text{thế kỉ}\,12\,\text{năm} = \dotfill\,\text{năm}$\par\vspace{1.4mm}
  \noindent
  d) $4\,\text{phút}\,35\,\text{giây} = \dotfill\,\text{giây}$\par
}

\vspace{0.8mm}

% BÀI 4: BÀI TOÁN GIẢI (11 DÒNG KẺ)
\khungBaiTap{%
  \textbf{Bài 4 (1,5 điểm).} Một cửa hàng nhập về 12{,}5 kg gạo. Mỗi ki-lô-gam gạo giá 18 000 đồng. Cửa hàng bán lại với giá 22 500 đồng mỗi ki-lô-gam. Hỏi sau khi bán hết số gạo đó, cửa hàng lãi bao nhiêu tiền?
  \par\vspace{0.3mm}%
  \noindent\textit{Bài giải:}
  \dongke{11}
}

\vspace{0.8mm}

% BÀI 5: TOÁN NÂNG CAO (8 DÒNG KẺ)
\khungBaiTap{%
  \textbf{Bài 5 (1,0 điểm).} Tính bằng cách thuận tiện nhất:\par\vspace{0.6mm}\centerline{$K = 3{,}5 \times 4{,}2 + 3{,}5 \times 5{,}8$}
  \par\vspace{0.3mm}%
  \noindent\textit{Bài làm:}
  \dongke{8}
}

\newpage


% =============================================================================
% ========================= MÃ ĐỀ: 402 ========================================
% =============================================================================

% -------------------- TRANG 1: MẶT TRƯỚC --------------------
\noindent
\begin{tabularx}{\linewidth}{|X|c|}
\hline
\textbf{PHIẾU BÀI TẬP CUỐI TUẦN -- TOÁN LỚP 4 (TUẦN 11)} & \textbf{MÃ ĐỀ: 402} \\
\textit{Trung tâm Toán Thầy Lĩnh -- Năm học: 2026 - 2027} & \textit{Thời gian: 45 phút} \\
\hline
\end{tabularx}

\vspace{1mm}

\noindent
\begin{tabularx}{\linewidth}{|X|p{3.5cm}|}
\hline
\rule[-2mm]{0pt}{7.5mm}\textbf{Họ và tên:} \dotfill & \textbf{Lớp:} \makebox[2cm]{\dotfill} \\
\hline
\end{tabularx}

\vspace{1mm}

\noindent
\begin{tabularx}{\linewidth}{|p{3cm}|X|}
\hline
\centering\textbf{ĐIỂM SỐ} & \centering\textbf{LỜI NHẬN XÉT CỦA THẦY CÔ} \tabularnewline
\hline
\rule{0pt}{1.1cm} & \tabularnewline
\hline
\end{tabularx}

\vspace{1mm}

\noindent\textbf{BẢNG ĐIỀN ĐÁP ÁN TRẮC NGHIỆM (Mỗi câu đúng 0,25 điểm):}
\vspace{0.6mm}

\begin{tabularx}{\linewidth}{|c|*{12}{>{\centering\arraybackslash}X|}}
\hline
\textbf{Câu} & \textbf{1} & \textbf{2} & \textbf{3} & \textbf{4} & \textbf{5} & \textbf{6} & \textbf{7} & \textbf{8} & \textbf{9} & \textbf{10} & \textbf{11} & \textbf{12} \\
\hline
\textbf{Đáp án} & & & & & & & & & & & & \\
\hline
\end{tabularx}

\vspace{1.2mm}

\noindent\textbf{\large PHẦN I. TRẮC NGHIỆM KHÁCH QUAN (3,0 điểm)}

\cauTN{1}{Kết quả của phép tính $348\,520 + 264\,395$ là:}{$612\,915$}{$612\,815$}{$602\,915$}{$613\,915$}
\cauTN{2}{Khi đổi chỗ các số hạng trong một tổng thì tổng đó:}{Thay đổi}{Tăng lên}{Không thay đổi}{Giảm đi}
\cauTN{3}{Giá trị của biểu thức $m + n + p$ với $m = 24\,000, n = 36\,000, p = 40\,000$ là:}{$90\,000$}{$100\,000$}{$110\,000$}{$80\,000$}
\cauTN{4}{Số thích hợp điền vào chỗ chấm $4\,\text{tấn}\,450\,\text{kg} = \dots\,\text{kg}$ là:}{$445$}{$4\,450$}{$40\,450$}{$44\,500$}
\cauTN{5}{Số thích hợp điền vào chỗ chấm $8\,\text{m}^2\,36\,\text{dm}^2 = \dots\,\text{dm}^2$ là:}{$836$}{$8\,036$}{$8\,360$}{$86$}
\cauTN{6}{Khoảng thời gian $4\,\text{phút}\,15\,\text{giây}$ bằng bao nhiêu giây?}{$245\,\text{giây}$}{$255\,\text{giây}$}{$265\,\text{giây}$}{$275\,\text{giây}$}
\cauTN{7}{Tìm số hạng chưa biết: $x + 42\,350 = 96\,820$:}{$x = 54\,470$}{$x = 54\,570$}{$x = 139\,170$}{$x = 54\,370$}
\cauTN{8}{Điền dấu thích hợp: $56\,000 + 27\,000 \;\dots\; 60\,000 + 24\,000$}{$>$}{$<$}{$=$}{Không so sánh được}
\cauTN{9}{Số bé nhất có $6$ chữ số khác nhau là:}{$100\,000$}{$102\,345$}{$123\,456$}{$102\,354$}
\cauTN{10}{Làm tròn số $583\,920$ đến hàng trăm nghìn ta được số:}{$500\,000$}{$580\,000$}{$600\,000$}{$590\,000$}
\cauTNhaiDong{11}{Một mảnh đất hình chữ nhật có chiều dài $50\,\text{m}$, chiều rộng $30\,\text{m}$. Chu vi mảnh đất là:}{$80\,\text{m}$}{$160\,\text{m}$}{$150\,\text{m}$}{$1\,500\,\text{m}^2$}
\cauTNhaiDong{12}{Tháng 1 một nhà máy sản xuất $1\,400$ sản phẩm, tháng 2 sản xuất $1\,600$ sản phẩm. Cả hai tháng nhà máy sản xuất được:}{$2\,800$}{$3\,000$}{$3\,200$}{$3\,500$}

\vspace{0.4mm}
\noindent\textbf{\large PHẦN II. TỰ LUẬN (7,0 điểm)}

\noindent\textbf{Bài 1 (2,0 điểm). Đặt tính rồi tính:}
\vspace{0.4mm}

\noindent
\begin{tabularx}{\linewidth}{|X|X|X|X|}
\hline
\oDatTinh[4]{a)}{$426\,735 + 185\,468$} & 
\oDatTinh[4]{b)}{$374\,829 + 258\,394$} & 
\oDatTinh[4]{c)}{$619\,482 + 84\,539$} & 
\oDatTinh[4]{d)}{$538\,246 + 461\,754$} \tabularnewline
\hline
\end{tabularx}

\newpage
% -------------------- TRANG 2: MẶT SAU --------------------
\noindent
\begin{tabularx}{\linewidth}{|X|c|}
\hline
\textbf{PHẦN TỰ LUẬN (TIẾP THEO) -- PHIẾU BÀI TẬP TUẦN 11} & \textbf{MÃ ĐỀ: 402} \\
\hline
\end{tabularx}

\vspace{0.8mm}

% BÀI 2: TÌM X (4 Ý)
\khungBaiTap{%
  \textbf{Bài 2 (1,5 điểm). Tìm $x$:}\par\vspace{0.6mm}
  \noindent
  \begin{minipage}[t]{0.485\linewidth}
    \textbf{a)} $x - 38\,450 = 84\,260$\par
    \dongke{3}
    \vspace{0.4mm}
    \textbf{c)} $85\,300 - x = 39\,460$\par
    \dongke{3}
  \end{minipage}\hfill
  \begin{minipage}[t]{0.485\linewidth}
    \textbf{b)} $x + 45\,380 = 98\,620$\par
    \dongke{3}
    \vspace{0.4mm}
    \textbf{d)} $x : 7 = 16\,342$\par
    \dongke{3}
  \end{minipage}
}

\vspace{0.8mm}

% BÀI 3: ĐỔI ĐƠN VỊ (4 Ý - 1 CỘT TỪ TRÊN XUỐNG)
\khungBaiTap{%
  \textbf{Bài 3 (1,0 điểm). Viết số thích hợp vào chỗ chấm:}\par\vspace{0.8mm}
  \noindent
  a) $7\,\text{tấn}\,420\,\text{kg} = \dotfill\,\text{kg}$\par\vspace{1.4mm}
  \noindent
  b) $9\,\text{m}^2\,18\,\text{dm}^2 = \dotfill\,\text{dm}^2$\par\vspace{1.4mm}
  \noindent
  c) $6\,\text{thế kỉ}\,25\,\text{năm} = \dotfill\,\text{năm}$\par\vspace{1.4mm}
  \noindent
  d) $3\,\text{phút}\,50\,\text{giây} = \dotfill\,\text{giây}$\par
}

\vspace{0.8mm}

% BÀI 4: BÀI TOÁN GIẢI (11 DÒNG KẺ)
\khungBaiTap{%
  \textbf{Bài 4 (1,5 điểm).} Một ô tô cứ đi 100 km thì tiêu thụ hết 8{,}5 lít xăng. Hỏi ô tô đó đi quãng đường 250 km thì tiêu thụ hết bao nhiêu lít xăng? Biết giá xăng là 24 500 đồng một lít, tính số tiền xăng phải trả.
  \par\vspace{0.3mm}%
  \noindent\textit{Bài giải:}
  \dongke{11}
}

\vspace{0.8mm}

% BÀI 5: TOÁN NÂNG CAO (8 DÒNG KẺ)
\khungBaiTap{%
  \textbf{Bài 5 (1,0 điểm).} Tính bằng cách thuận tiện nhất:\par\vspace{0.6mm}\centerline{$K = 2{,}4 \times 6{,}7 + 2{,}4 \times 3{,}3$}
  \par\vspace{0.3mm}%
  \noindent\textit{Bài làm:}
  \dongke{8}
}

\newpage


% =============================================================================
% ========================= MÃ ĐỀ: 403 ========================================
% =============================================================================

% -------------------- TRANG 1: MẶT TRƯỚC --------------------
\noindent
\begin{tabularx}{\linewidth}{|X|c|}
\hline
\textbf{PHIẾU BÀI TẬP CUỐI TUẦN -- TOÁN LỚP 4 (TUẦN 11)} & \textbf{MÃ ĐỀ: 403} \\
\textit{Trung tâm Toán Thầy Lĩnh -- Năm học: 2026 - 2027} & \textit{Thời gian: 45 phút} \\
\hline
\end{tabularx}

\vspace{1mm}

\noindent
\begin{tabularx}{\linewidth}{|X|p{3.5cm}|}
\hline
\rule[-2mm]{0pt}{7.5mm}\textbf{Họ và tên:} \dotfill & \textbf{Lớp:} \makebox[2cm]{\dotfill} \\
\hline
\end{tabularx}

\vspace{1mm}

\noindent
\begin{tabularx}{\linewidth}{|p{3cm}|X|}
\hline
\centering\textbf{ĐIỂM SỐ} & \centering\textbf{LỜI NHẬN XÉT CỦA THẦY CÔ} \tabularnewline
\hline
\rule{0pt}{1.1cm} & \tabularnewline
\hline
\end{tabularx}

\vspace{1mm}

\noindent\textbf{BẢNG ĐIỀN ĐÁP ÁN TRẮC NGHIỆM (Mỗi câu đúng 0,25 điểm):}
\vspace{0.6mm}

\begin{tabularx}{\linewidth}{|c|*{12}{>{\centering\arraybackslash}X|}}
\hline
\textbf{Câu} & \textbf{1} & \textbf{2} & \textbf{3} & \textbf{4} & \textbf{5} & \textbf{6} & \textbf{7} & \textbf{8} & \textbf{9} & \textbf{10} & \textbf{11} & \textbf{12} \\
\hline
\textbf{Đáp án} & & & & & & & & & & & & \\
\hline
\end{tabularx}

\vspace{1.2mm}

\noindent\textbf{\large PHẦN I. TRẮC NGHIỆM KHÁCH QUAN (3,0 điểm)}

\cauTN{1}{Kết quả của phép tính $452\,630 + 179\,485$ là:}{$632\,115$}{$631\,115$}{$622\,115$}{$632\,015$}
\cauTN{2}{Điền vào chỗ chấm: $m + n = \dots + m$}{$n$}{$m$}{$0$}{$1$}
\cauTN{3}{Giá trị của biểu thức $a + b + c$ với $a = 35\,000, b = 45\,000, c = 20\,000$ là:}{$90\,000$}{$100\,000$}{$110\,000$}{$120\,000$}
\cauTN{4}{Số thích hợp điền vào chỗ chấm $6\,\text{tấn}\,180\,\text{kg} = \dots\,\text{kg}$ là:}{$618$}{$6\,180$}{$60\,180$}{$61\,800$}
\cauTN{5}{Số thích hợp điền vào chỗ chấm $5\,\text{m}^2\,62\,\text{dm}^2 = \dots\,\text{dm}^2$ là:}{$562$}{$5\,062$}{$5\,620$}{$58$}
\cauTN{6}{Khoảng thời gian $5\,\text{phút}\,20\,\text{giây}$ bằng bao nhiêu giây?}{$300\,\text{giây}$}{$320\,\text{giây}$}{$340\,\text{giây}$}{$360\,\text{giây}$}
\cauTN{7}{Tìm số hạng chưa biết: $x + 28\,640 = 75\,920$:}{$x = 47\,280$}{$x = 47\,380$}{$x = 104\,560$}{$x = 46\,280$}
\cauTN{8}{Điền dấu thích hợp: $68\,000 + 14\,000 \;\dots\; 50\,000 + 32\,000$}{$>$}{$<$}{$=$}{Không so sánh được}
\cauTN{9}{Số chẵn lớn nhất có $6$ chữ số là:}{$999\,998$}{$999\,999$}{$987\,654$}{$999\,990$}
\cauTN{10}{Làm tròn số $635\,480$ đến hàng trăm nghìn ta được số:}{$600\,000$}{$700\,000$}{$640\,000$}{$630\,000$}
\cauTNhaiDong{11}{Một sân trường hình chữ nhật có chu vi $180\,\text{m}$. Nửa chu vi sân trường là:}{$45\,\text{m}$}{$90\,\text{m}$}{$60\,\text{m}$}{$360\,\text{m}$}
\cauTNhaiDong{12}{Một đội công nhân tuần thứ nhất trải được $1\,250\,\text{m}$ đường, tuần thứ hai trải được $1\,750\,\text{m}$. Cả hai tuần đội trải được số ki-lô-mét đường là:}{$2\,\text{km}$}{$3\,\text{km}$}{$4\,\text{km}$}{$30\,\text{km}$}

\vspace{0.4mm}
\noindent\textbf{\large PHẦN II. TỰ LUẬN (7,0 điểm)}

\noindent\textbf{Bài 1 (2,0 điểm). Đặt tính rồi tính:}
\vspace{0.4mm}

\noindent
\begin{tabularx}{\linewidth}{|X|X|X|X|}
\hline
\oDatTinh[4]{a)}{$518\,462 + 237\,589$} & 
\oDatTinh[4]{b)}{$462\,395 + 178\,648$} & 
\oDatTinh[4]{c)}{$724\,536 + 68\,475$} & 
\oDatTinh[4]{d)}{$645\,128 + 354\,872$} \tabularnewline
\hline
\end{tabularx}

\newpage
% -------------------- TRANG 2: MẶT SAU --------------------
\noindent
\begin{tabularx}{\linewidth}{|X|c|}
\hline
\textbf{PHẦN TỰ LUẬN (TIẾP THEO) -- PHIẾU BÀI TẬP TUẦN 11} & \textbf{MÃ ĐỀ: 403} \\
\hline
\end{tabularx}

\vspace{0.8mm}

% BÀI 2: TÌM X (4 Ý)
\khungBaiTap{%
  \textbf{Bài 2 (1,5 điểm). Tìm $x$:}\par\vspace{0.6mm}
  \noindent
  \begin{minipage}[t]{0.485\linewidth}
    \textbf{a)} $x - 52\,360 = 69\,480$\par
    \dongke{3}
    \vspace{0.4mm}
    \textbf{c)} $92\,400 - x = 45\,680$\par
    \dongke{3}
  \end{minipage}\hfill
  \begin{minipage}[t]{0.485\linewidth}
    \textbf{b)} $x + 54\,270 = 89\,630$\par
    \dongke{3}
    \vspace{0.4mm}
    \textbf{d)} $x : 6 = 18\,245$\par
    \dongke{3}
  \end{minipage}
}

\vspace{0.8mm}

% BÀI 3: ĐỔI ĐƠN VỊ (4 Ý - 1 CỘT TỪ TRÊN XUỐNG)
\khungBaiTap{%
  \textbf{Bài 3 (1,0 điểm). Viết số thích hợp vào chỗ chấm:}\par\vspace{0.8mm}
  \noindent
  a) $8\,\text{tấn}\,260\,\text{kg} = \dotfill\,\text{kg}$\par\vspace{1.4mm}
  \noindent
  b) $6\,\text{m}^2\,54\,\text{dm}^2 = \dotfill\,\text{dm}^2$\par\vspace{1.4mm}
  \noindent
  c) $4\,\text{thế kỉ}\,30\,\text{năm} = \dotfill\,\text{năm}$\par\vspace{1.4mm}
  \noindent
  d) $5\,\text{phút}\,12\,\text{giây} = \dotfill\,\text{giây}$\par
}

\vspace{0.8mm}

% BÀI 4: BÀI TOÁN GIẢI (11 DÒNG KẺ)
\khungBaiTap{%
  \textbf{Bài 4 (1,5 điểm).} Một cửa hàng văn phòng phẩm bán 15{,}5 tá bút chì. Mỗi tá bút chì giá 78 000 đồng. Cửa hàng giảm giá 5 000 đồng cho mỗi tá. Hỏi cửa hàng thu được bao nhiêu tiền khi bán hết số bút chì đó?
  \par\vspace{0.3mm}%
  \noindent\textit{Bài giải:}
  \dongke{11}
}

\vspace{0.8mm}

% BÀI 5: TOÁN NÂNG CAO (8 DÒNG KẺ)
\khungBaiTap{%
  \textbf{Bài 5 (1,0 điểm).} Tính bằng cách thuận tiện nhất:\par\vspace{0.6mm}\centerline{$K = 0{,}8 \times 12{,}5 + 0{,}8 \times 7{,}5$}
  \par\vspace{0.3mm}%
  \noindent\textit{Bài làm:}
  \dongke{8}
}

\newpage


% =============================================================================
% ========================= MÃ ĐỀ: 404 ========================================
% =============================================================================

% -------------------- TRANG 1: MẶT TRƯỚC --------------------
\noindent
\begin{tabularx}{\linewidth}{|X|c|}
\hline
\textbf{PHIẾU BÀI TẬP CUỐI TUẦN -- TOÁN LỚP 4 (TUẦN 11)} & \textbf{MÃ ĐỀ: 404} \\
\textit{Trung tâm Toán Thầy Lĩnh -- Năm học: 2026 - 2027} & \textit{Thời gian: 45 phút} \\
\hline
\end{tabularx}

\vspace{1mm}

\noindent
\begin{tabularx}{\linewidth}{|X|p{3.5cm}|}
\hline
\rule[-2mm]{0pt}{7.5mm}\textbf{Họ và tên:} \dotfill & \textbf{Lớp:} \makebox[2cm]{\dotfill} \\
\hline
\end{tabularx}

\vspace{1mm}

\noindent
\begin{tabularx}{\linewidth}{|p{3cm}|X|}
\hline
\centering\textbf{ĐIỂM SỐ} & \centering\textbf{LỜI NHẬN XÉT CỦA THẦY CÔ} \tabularnewline
\hline
\rule{0pt}{1.1cm} & \tabularnewline
\hline
\end{tabularx}

\vspace{1mm}

\noindent\textbf{BẢNG ĐIỀN ĐÁP ÁN TRẮC NGHIỆM (Mỗi câu đúng 0,25 điểm):}
\vspace{0.6mm}

\begin{tabularx}{\linewidth}{|c|*{12}{>{\centering\arraybackslash}X|}}
\hline
\textbf{Câu} & \textbf{1} & \textbf{2} & \textbf{3} & \textbf{4} & \textbf{5} & \textbf{6} & \textbf{7} & \textbf{8} & \textbf{9} & \textbf{10} & \textbf{11} & \textbf{12} \\
\hline
\textbf{Đáp án} & & & & & & & & & & & & \\
\hline
\end{tabularx}

\vspace{1.2mm}

\noindent\textbf{\large PHẦN I. TRẮC NGHIỆM KHÁCH QUAN (3,0 điểm)}

\cauTN{1}{Kết quả của phép tính $536\,420 + 285\,694$ là:}{$822\,114$}{$821\,114$}{$812\,114$}{$822\,014$}
\cauTN{2}{Phép tính nào dưới đây có kết quả bằng $15\,000 + 25\,000$?}{$25\,000 + 15\,000$}{$25\,000 - 15\,000$}{$15\,000 \times 2$}{$40\,000 + 10\,000$}
\cauTN{3}{Giá trị của biểu thức $x + y + z$ với $x = 42\,000, y = 18\,000, z = 50\,000$ là:}{$100\,000$}{$110\,000$}{$120\,000$}{$90\,000$}
\cauTN{4}{Số thích hợp điền vào chỗ chấm $7\,\text{tấn}\,60\,\text{kg} = \dots\,\text{kg}$ là:}{$760$}{$7\,060$}{$7\,600$}{$76\,000$}
\cauTN{5}{Số thích hợp điền vào chỗ chấm $4\,\text{m}^2\,85\,\text{dm}^2 = \dots\,\text{dm}^2$ là:}{$485$}{$4\,085$}{$4\,850$}{$49$}
\cauTN{6}{Khoảng thời gian $2\,\text{phút}\,55\,\text{giây}$ bằng bao nhiêu giây?}{$165\,\text{giây}$}{$175\,\text{giây}$}{$185\,\text{giây}$}{$195\,\text{giây}$}
\cauTN{7}{Tìm số hạng chưa biết: $x + 37\,520 = 84\,690$:}{$x = 47\,170$}{$x = 47\,270$}{$x = 122\,210$}{$x = 46\,170$}
\cauTN{8}{Điền dấu thích hợp: $72\,000 + 18\,000 \;\dots\; 65\,000 + 26\,000$}{$>$}{$<$}{$=$}{Không so sánh được}
\cauTN{9}{Số lẻ bé nhất có $6$ chữ số là:}{$100\,000$}{$100\,001$}{$100\,003$}{$102\,345$}
\cauTN{10}{Làm tròn số $749\,200$ đến hàng trăm nghìn ta được số:}{$700\,000$}{$750\,000$}{$800\,000$}{$740\,000$}
\cauTNhaiDong{11}{Một sân bóng đá hình chữ nhật có chu vi $320\,\text{m}$. Nửa chu vi sân bóng là:}{$80\,\text{m}$}{$160\,\text{m}$}{$100\,\text{m}$}{$640\,\text{m}$}
\cauTNhaiDong{12}{Một cửa hàng bán được $2\,500\,\text{m}$ vải hoa và $1\,500\,\text{m}$ vải trắng. Tổng số ki-lô-mét vải bán được là:}{$3\,\text{km}$}{$4\,\text{km}$}{$5\,\text{km}$}{$40\,\text{km}$}

\vspace{0.4mm}
\noindent\textbf{\large PHẦN II. TỰ LUẬN (7,0 điểm)}

\noindent\textbf{Bài 1 (2,0 điểm). Đặt tính rồi tính:}
\vspace{0.4mm}

\noindent
\begin{tabularx}{\linewidth}{|X|X|X|X|}
\hline
\oDatTinh[4]{a)}{$364\,752 + 287\,369$} & 
\oDatTinh[4]{b)}{$528\,416 + 194\,785$} & 
\oDatTinh[4]{c)}{$635\,284 + 86\,927$} & 
\oDatTinh[4]{d)}{$724\,350 + 275\,650$} \tabularnewline
\hline
\end{tabularx}

\newpage
% -------------------- TRANG 2: MẶT SAU --------------------
\noindent
\begin{tabularx}{\linewidth}{|X|c|}
\hline
\textbf{PHẦN TỰ LUẬN (TIẾP THEO) -- PHIẾU BÀI TẬP TUẦN 11} & \textbf{MÃ ĐỀ: 404} \\
\hline
\end{tabularx}

\vspace{0.8mm}

% BÀI 2: TÌM X (4 Ý)
\khungBaiTap{%
  \textbf{Bài 2 (1,5 điểm). Tìm $x$:}\par\vspace{0.6mm}
  \noindent
  \begin{minipage}[t]{0.485\linewidth}
    \textbf{a)} $x - 46\,290 = 75\,830$\par
    \dongke{3}
    \vspace{0.4mm}
    \textbf{c)} $86\,200 - x = 37\,490$\par
    \dongke{3}
  \end{minipage}\hfill
  \begin{minipage}[t]{0.485\linewidth}
    \textbf{b)} $x + 38\,460 = 82\,150$\par
    \dongke{3}
    \vspace{0.4mm}
    \textbf{d)} $x : 9 = 15\,234$\par
    \dongke{3}
  \end{minipage}
}

\vspace{0.8mm}

% BÀI 3: ĐỔI ĐƠN VỊ (4 Ý - 1 CỘT TỪ TRÊN XUỐNG)
\khungBaiTap{%
  \textbf{Bài 3 (1,0 điểm). Viết số thích hợp vào chỗ chấm:}\par\vspace{0.8mm}
  \noindent
  a) $9\,\text{tấn}\,85\,\text{kg} = \dotfill\,\text{kg}$\par\vspace{1.4mm}
  \noindent
  b) $7\,\text{m}^2\,92\,\text{dm}^2 = \dotfill\,\text{dm}^2$\par\vspace{1.4mm}
  \noindent
  c) $8\,\text{thế kỉ}\,40\,\text{năm} = \dotfill\,\text{năm}$\par\vspace{1.4mm}
  \noindent
  d) $6\,\text{phút}\,10\,\text{giây} = \dotfill\,\text{giây}$\par
}

\vspace{0.8mm}

% BÀI 4: BÀI TOÁN GIẢI (11 DÒNG KẺ)
\khungBaiTap{%
  \textbf{Bài 4 (1,5 điểm).} Một xe máy cứ đi 100 km thì tiêu thụ hết 2{,}5 lít xăng. Hỏi xe máy đó đi quãng đường 180 km thì tiêu thụ hết bao nhiêu lít xăng? Nếu giá xăng là 23 000 đồng một lít thì đi 180 km hết bao nhiêu tiền xăng?
  \par\vspace{0.3mm}%
  \noindent\textit{Bài giải:}
  \dongke{11}
}

\vspace{0.8mm}

% BÀI 5: TOÁN NÂNG CAO (8 DÒNG KẺ)
\khungBaiTap{%
  \textbf{Bài 5 (1,0 điểm).} Tính bằng cách thuận tiện nhất:\par\vspace{0.6mm}\centerline{$K = 4{,}5 \times 3{,}6 + 4{,}5 \times 6{,}4$}
  \par\vspace{0.3mm}%
  \noindent\textit{Bài làm:}
  \dongke{8}
}

\newpage


% =============================================================================
% ========================= MÃ ĐỀ: 405 ========================================
% =============================================================================

% -------------------- TRANG 1: MẶT TRƯỚC --------------------
\noindent
\begin{tabularx}{\linewidth}{|X|c|}
\hline
\textbf{PHIẾU BÀI TẬP CUỐI TUẦN -- TOÁN LỚP 4 (TUẦN 11)} & \textbf{MÃ ĐỀ: 405} \\
\textit{Trung tâm Toán Thầy Lĩnh -- Năm học: 2026 - 2027} & \textit{Thời gian: 45 phút} \\
\hline
\end{tabularx}

\vspace{1mm}

\noindent
\begin{tabularx}{\linewidth}{|X|p{3.5cm}|}
\hline
\rule[-2mm]{0pt}{7.5mm}\textbf{Họ và tên:} \dotfill & \textbf{Lớp:} \makebox[2cm]{\dotfill} \\
\hline
\end{tabularx}

\vspace{1mm}

\noindent
\begin{tabularx}{\linewidth}{|p{3cm}|X|}
\hline
\centering\textbf{ĐIỂM SỐ} & \centering\textbf{LỜI NHẬN XÉT CỦA THẦY CÔ} \tabularnewline
\hline
\rule{0pt}{1.1cm} & \tabularnewline
\hline
\end{tabularx}

\vspace{1mm}

\noindent\textbf{BẢNG ĐIỀN ĐÁP ÁN TRẮC NGHIỆM (Mỗi câu đúng 0,25 điểm):}
\vspace{0.6mm}

\begin{tabularx}{\linewidth}{|c|*{12}{>{\centering\arraybackslash}X|}}
\hline
\textbf{Câu} & \textbf{1} & \textbf{2} & \textbf{3} & \textbf{4} & \textbf{5} & \textbf{6} & \textbf{7} & \textbf{8} & \textbf{9} & \textbf{10} & \textbf{11} & \textbf{12} \\
\hline
\textbf{Đáp án} & & & & & & & & & & & & \\
\hline
\end{tabularx}

\vspace{1.2mm}

\noindent\textbf{\large PHẦN I. TRẮC NGHIỆM KHÁCH QUAN (3,0 điểm)}

\cauTN{1}{Kết quả của phép tính $618\,540 + 293\,685$ là:}{$912\,225$}{$911\,225$}{$902\,225$}{$912\,125$}
\cauTN{2}{Tính chất giao hoán của phép cộng cho biết: $a + b = \dots$}{$b + a$}{$a - b$}{$b \times a$}{$a + 0$}
\cauTN{3}{Giá trị của biểu thức $a + b + c$ với $a = 55\,000, b = 25\,000, c = 40\,000$ là:}{$110\,000$}{$120\,000$}{$130\,000$}{$100\,000$}
\cauTN{4}{Số thích hợp điền vào chỗ chấm $8\,\text{tấn}\,320\,\text{kg} = \dots\,\text{kg}$ là:}{$832$}{$8\,320$}{$80\,320$}{$83\,200$}
\cauTN{5}{Số thích hợp điền vào chỗ chấm $6\,\text{m}^2\,19\,\text{dm}^2 = \dots\,\text{dm}^2$ là:}{$619$}{$6\,019$}{$6\,190$}{$63$}
\cauTN{6}{Khoảng thời gian $3\,\text{phút}\,40\,\text{giây}$ bằng bao nhiêu giây?}{$200\,\text{giây}$}{$220\,\text{giây}$}{$240\,\text{giây}$}{$260\,\text{giây}$}
\cauTN{7}{Tìm số hạng chưa biết: $x + 46\,380 = 91\,520$:}{$x = 45\,140$}{$x = 45\,240$}{$x = 137\,900$}{$x = 44\,140$}
\cauTN{8}{Điền dấu thích hợp: $83\,000 + 17\,000 \;\dots\; 75\,000 + 25\,000$}{$>$}{$<$}{$=$}{Không so sánh được}
\cauTN{9}{Số tròn chục nghìn lớn nhất có $6$ chữ số là:}{$990\,000$}{$999\,990$}{$980\,000$}{$900\,000$}
\cauTN{10}{Làm tròn số $852\,600$ đến hàng trăm nghìn ta được số:}{$800\,000$}{$850\,000$}{$900\,000$}{$860\,000$}
\cauTNhaiDong{11}{Chu vi một mảnh vườn hình chữ nhật là $240\,\text{m}$. Nửa chu vi mảnh vườn là:}{$60\,\text{m}$}{$120\,\text{m}$}{$180\,\text{m}$}{$480\,\text{m}$}
\cauTNhaiDong{12}{Bác An thu hoạch vụ mùa thứ nhất được $2\,200\,\text{kg}$ thóc, vụ thứ hai được $1\,800\,\text{kg}$. Tổng sản lượng cả hai vụ tính theo tấn là:}{$3\,\text{tấn}$}{$4\,\text{tấn}$}{$5\,\text{tấn}$}{$40\,\text{tấn}$}

\vspace{0.4mm}
\noindent\textbf{\large PHẦN II. TỰ LUẬN (7,0 điểm)}

\noindent\textbf{Bài 1 (2,0 điểm). Đặt tính rồi tính:}
\vspace{0.4mm}

\noindent
\begin{tabularx}{\linewidth}{|X|X|X|X|}
\hline
\oDatTinh[4]{a)}{$452\,836 + 269\,478$} & 
\oDatTinh[4]{b)}{$384\,629 + 257\,894$} & 
\oDatTinh[4]{c)}{$548\,392 + 95\,746$} & 
\oDatTinh[4]{d)}{$483\,165 + 516\,835$} \tabularnewline
\hline
\end{tabularx}

\newpage
% -------------------- TRANG 2: MẶT SAU --------------------
\noindent
\begin{tabularx}{\linewidth}{|X|c|}
\hline
\textbf{PHẦN TỰ LUẬN (TIẾP THEO) -- PHIẾU BÀI TẬP TUẦN 11} & \textbf{MÃ ĐỀ: 405} \\
\hline
\end{tabularx}

\vspace{0.8mm}

% BÀI 2: TÌM X (4 Ý)
\khungBaiTap{%
  \textbf{Bài 2 (1,5 điểm). Tìm $x$:}\par\vspace{0.6mm}
  \noindent
  \begin{minipage}[t]{0.485\linewidth}
    \textbf{a)} $x - 58\,340 = 67\,890$\par
    \dongke{3}
    \vspace{0.4mm}
    \textbf{c)} $95\,300 - x = 48\,620$\par
    \dongke{3}
  \end{minipage}\hfill
  \begin{minipage}[t]{0.485\linewidth}
    \textbf{b)} $x + 47\,620 = 93\,480$\par
    \dongke{3}
    \vspace{0.4mm}
    \textbf{d)} $x : 8 = 17\,345$\par
    \dongke{3}
  \end{minipage}
}

\vspace{0.8mm}

% BÀI 3: ĐỔI ĐƠN VỊ (4 Ý - 1 CỘT TỪ TRÊN XUỐNG)
\khungBaiTap{%
  \textbf{Bài 3 (1,0 điểm). Viết số thích hợp vào chỗ chấm:}\par\vspace{0.8mm}
  \noindent
  a) $5\,\text{tấn}\,720\,\text{kg} = \dotfill\,\text{kg}$\par\vspace{1.4mm}
  \noindent
  b) $8\,\text{m}^2\,63\,\text{dm}^2 = \dotfill\,\text{dm}^2$\par\vspace{1.4mm}
  \noindent
  c) $7\,\text{thế kỉ}\,15\,\text{năm} = \dotfill\,\text{năm}$\par\vspace{1.4mm}
  \noindent
  d) $4\,\text{phút}\,45\,\text{giây} = \dotfill\,\text{giây}$\par
}

\vspace{0.8mm}

% BÀI 4: BÀI TOÁN GIẢI (11 DÒNG KẺ)
\khungBaiTap{%
  \textbf{Bài 4 (1,5 điểm).} Một cửa hàng nhập về 25{,}5 kg đường. Mỗi ki-lô-gam đường giá 16 500 đồng. Cửa hàng đã bán được $\frac{2}{5}$ số đường với giá 21 000 đồng mỗi ki-lô-gam. Hỏi cửa hàng còn lại bao nhiêu ki-lô-gam đường và đã thu được bao nhiêu tiền từ số đường đã bán?
  \par\vspace{0.3mm}%
  \noindent\textit{Bài giải:}
  \dongke{11}
}

\vspace{0.8mm}

% BÀI 5: TOÁN NÂNG CAO (8 DÒNG KẺ)
\khungBaiTap{%
  \textbf{Bài 5 (1,0 điểm).} Tính bằng cách thuận tiện nhất:\par\vspace{0.6mm}\centerline{$K = 1{,}25 \times 7{,}8 + 1{,}25 \times 2{,}2$}
  \par\vspace{0.3mm}%
  \noindent\textit{Bài làm:}
  \dongke{8}
}

\newpage


% =============================================================================
% ========================= MÃ ĐỀ: 406 ========================================
% =============================================================================

% -------------------- TRANG 1: MẶT TRƯỚC --------------------
\noindent
\begin{tabularx}{\linewidth}{|X|c|}
\hline
\textbf{PHIẾU BÀI TẬP CUỐI TUẦN -- TOÁN LỚP 4 (TUẦN 11)} & \textbf{MÃ ĐỀ: 406} \\
\textit{Trung tâm Toán Thầy Lĩnh -- Năm học: 2026 - 2027} & \textit{Thời gian: 45 phút} \\
\hline
\end{tabularx}

\vspace{1mm}

\noindent
\begin{tabularx}{\linewidth}{|X|p{3.5cm}|}
\hline
\rule[-2mm]{0pt}{7.5mm}\textbf{Họ và tên:} \dotfill & \textbf{Lớp:} \makebox[2cm]{\dotfill} \\
\hline
\end{tabularx}

\vspace{1mm}

\noindent
\begin{tabularx}{\linewidth}{|p{3cm}|X|}
\hline
\centering\textbf{ĐIỂM SỐ} & \centering\textbf{LỜI NHẬN XÉT CỦA THẦY CÔ} \tabularnewline
\hline
\rule{0pt}{1.1cm} & \tabularnewline
\hline
\end{tabularx}

\vspace{1mm}

\noindent\textbf{BẢNG ĐIỀN ĐÁP ÁN TRẮC NGHIỆM (Mỗi câu đúng 0,25 điểm):}
\vspace{0.6mm}

\begin{tabularx}{\linewidth}{|c|*{12}{>{\centering\arraybackslash}X|}}
\hline
\textbf{Câu} & \textbf{1} & \textbf{2} & \textbf{3} & \textbf{4} & \textbf{5} & \textbf{6} & \textbf{7} & \textbf{8} & \textbf{9} & \textbf{10} & \textbf{11} & \textbf{12} \\
\hline
\textbf{Đáp án} & & & & & & & & & & & & \\
\hline
\end{tabularx}

\vspace{1.2mm}

\noindent\textbf{\large PHẦN I. TRẮC NGHIỆM KHÁCH QUAN (3,0 điểm)}

\cauTN{1}{Kết quả của phép tính $724\,360 + 189\,845$ là:}{$914\,205$}{$913\,205$}{$904\,205$}{$914\,105$}
\cauTN{2}{Nếu $a + b = c$ thì $b + a = \dots$}{$c$}{$a$}{$b$}{$2c$}
\cauTN{3}{Giá trị của biểu thức $a + b + c$ với $a = 48\,000, b = 32\,000, c = 50\,000$ là:}{$120\,000$}{$130\,000$}{$140\,000$}{$110\,000$}
\cauTN{4}{Số thích hợp điền vào chỗ chấm $3\,\text{tấn}\,850\,\text{kg} = \dots\,\text{kg}$ là:}{$385$}{$3\,850$}{$38\,500$}{$30\,850$}
\cauTN{5}{Số thích hợp điền vào chỗ chấm $9\,\text{m}^2\,28\,\text{dm}^2 = \dots\,\text{dm}^2$ là:}{$928$}{$9\,028$}{$9\,280$}{$95$}
\cauTN{6}{Khoảng thời gian $5\,\text{phút}\,45\,\text{giây}$ bằng bao nhiêu giây?}{$325\,\text{giây}$}{$345\,\text{giây}$}{$355\,\text{giây}$}{$365\,\text{giây}$}
\cauTN{7}{Tìm số hạng chưa biết: $x + 52\,470 = 87\,620$:}{$x = 35\,150$}{$x = 35\,250$}{$x = 140\,090$}{$x = 34\,150$}
\cauTN{8}{Điền dấu thích hợp: $64\,000 + 29\,000 \;\dots\; 70\,000 + 25\,000$}{$>$}{$<$}{$=$}{Không so sánh được}
\cauTN{9}{Số tròn trăm nghìn bé nhất có $6$ chữ số là:}{$100\,000$}{$200\,000$}{$100\,001$}{$900\,000$}
\cauTN{10}{Làm tròn số $938\,400$ đến hàng trăm nghìn ta được số:}{$900\,000$}{$940\,000$}{$1\,000\,000$}{$930\,000$}
\cauTNhaiDong{11}{Chu vi hình chữ nhật bằng $280\,\text{m}$. Nửa chu vi hình chữ nhật là:}{$70\,\text{m}$}{$140\,\text{m}$}{$210\,\text{m}$}{$560\,\text{m}$}
\cauTNhaiDong{12}{Một đoàn xe gồm $2$ xe tải, xe thứ nhất chở $3\,500\,\text{kg}$ hàng, xe thứ hai chở $2\,500\,\text{kg}$ hàng. Cả hai xe chở được bao nhiêu tấn hàng?}{$5\,\text{tấn}$}{$6\,\text{tấn}$}{$7\,\text{tấn}$}{$60\,\text{tấn}$}

\vspace{0.4mm}
\noindent\textbf{\large PHẦN II. TỰ LUẬN (7,0 điểm)}

\noindent\textbf{Bài 1 (2,0 điểm). Đặt tính rồi tính:}
\vspace{0.4mm}

\noindent
\begin{tabularx}{\linewidth}{|X|X|X|X|}
\hline
\oDatTinh[4]{a)}{$563\,724 + 278\,489$} & 
\oDatTinh[4]{b)}{$418\,635 + 294\,786$} & 
\oDatTinh[4]{c)}{$682\,419 + 78\,695$} & 
\oDatTinh[4]{d)}{$387\,452 + 612\,548$} \tabularnewline
\hline
\end{tabularx}

\newpage
% -------------------- TRANG 2: MẶT SAU --------------------
\noindent
\begin{tabularx}{\linewidth}{|X|c|}
\hline
\textbf{PHẦN TỰ LUẬN (TIẾP THEO) -- PHIẾU BÀI TẬP TUẦN 11} & \textbf{MÃ ĐỀ: 406} \\
\hline
\end{tabularx}

\vspace{0.8mm}

% BÀI 2: TÌM X (4 Ý)
\khungBaiTap{%
  \textbf{Bài 2 (1,5 điểm). Tìm $x$:}\par\vspace{0.6mm}
  \noindent
  \begin{minipage}[t]{0.485\linewidth}
    \textbf{a)} $x - 63\,480 = 74\,650$\par
    \dongke{3}
    \vspace{0.4mm}
    \textbf{c)} $97\,400 - x = 52\,680$\par
    \dongke{3}
  \end{minipage}\hfill
  \begin{minipage}[t]{0.485\linewidth}
    \textbf{b)} $x + 58\,340 = 95\,170$\par
    \dongke{3}
    \vspace{0.4mm}
    \textbf{d)} $x : 7 = 19\,425$\par
    \dongke{3}
  \end{minipage}
}

\vspace{0.8mm}

% BÀI 3: ĐỔI ĐƠN VỊ (4 Ý - 1 CỘT TỪ TRÊN XUỐNG)
\khungBaiTap{%
  \textbf{Bài 3 (1,0 điểm). Viết số thích hợp vào chỗ chấm:}\par\vspace{0.8mm}
  \noindent
  a) $4\,\text{tấn}\,640\,\text{kg} = \dotfill\,\text{kg}$\par\vspace{1.4mm}
  \noindent
  b) $5\,\text{m}^2\,73\,\text{dm}^2 = \dotfill\,\text{dm}^2$\par\vspace{1.4mm}
  \noindent
  c) $9\,\text{thế kỉ}\,50\,\text{năm} = \dotfill\,\text{năm}$\par\vspace{1.4mm}
  \noindent
  d) $6\,\text{phút}\,25\,\text{giây} = \dotfill\,\text{giây}$\par
}

\vspace{0.8mm}

% BÀI 4: BÀI TOÁN GIẢI (11 DÒNG KẺ)
\khungBaiTap{%
  \textbf{Bài 4 (1,5 điểm).} Một ô tô cứ đi 100 km thì tiêu thụ hết 9{,}2 lít xăng. Hỏi ô tô đó đi quãng đường 350 km thì tiêu thụ hết bao nhiêu lít xăng? Biết giá xăng là 25 000 đồng một lít, tính số tiền xăng phải trả.
  \par\vspace{0.3mm}%
  \noindent\textit{Bài giải:}
  \dongke{11}
}

\vspace{0.8mm}

% BÀI 5: TOÁN NÂNG CAO (8 DÒNG KẺ)
\khungBaiTap{%
  \textbf{Bài 5 (1,0 điểm).} Tính bằng cách thuận tiện nhất:\par\vspace{0.6mm}\centerline{$K = 6{,}4 \times 2{,}5 + 6{,}4 \times 7{,}5$}
  \par\vspace{0.3mm}%
  \noindent\textit{Bài làm:}
  \dongke{8}
}

\newpage


% =============================================================================
% ========================= MÃ ĐỀ: 407 ========================================
% =============================================================================

% -------------------- TRANG 1: MẶT TRƯỚC --------------------
\noindent
\begin{tabularx}{\linewidth}{|X|c|}
\hline
\textbf{PHIẾU BÀI TẬP CUỐI TUẦN -- TOÁN LỚP 4 (TUẦN 11)} & \textbf{MÃ ĐỀ: 407} \\
\textit{Trung tâm Toán Thầy Lĩnh -- Năm học: 2026 - 2027} & \textit{Thời gian: 45 phút} \\
\hline
\end{tabularx}

\vspace{1mm}

\noindent
\begin{tabularx}{\linewidth}{|X|p{3.5cm}|}
\hline
\rule[-2mm]{0pt}{7.5mm}\textbf{Họ và tên:} \dotfill & \textbf{Lớp:} \makebox[2cm]{\dotfill} \\
\hline
\end{tabularx}

\vspace{1mm}

\noindent
\begin{tabularx}{\linewidth}{|p{3cm}|X|}
\hline
\centering\textbf{ĐIỂM SỐ} & \centering\textbf{LỜI NHẬN XÉT CỦA THẦY CÔ} \tabularnewline
\hline
\rule{0pt}{1.1cm} & \tabularnewline
\hline
\end{tabularx}

\vspace{1mm}

\noindent\textbf{BẢNG ĐIỀN ĐÁP ÁN TRẮC NGHIỆM (Mỗi câu đúng 0,25 điểm):}
\vspace{0.6mm}

\begin{tabularx}{\linewidth}{|c|*{12}{>{\centering\arraybackslash}X|}}
\hline
\textbf{Câu} & \textbf{1} & \textbf{2} & \textbf{3} & \textbf{4} & \textbf{5} & \textbf{6} & \textbf{7} & \textbf{8} & \textbf{9} & \textbf{10} & \textbf{11} & \textbf{12} \\
\hline
\textbf{Đáp án} & & & & & & & & & & & & \\
\hline
\end{tabularx}

\vspace{1.2mm}

\noindent\textbf{\large PHẦN I. TRẮC NGHIỆM KHÁCH QUAN (3,0 điểm)}

\cauTN{1}{Kết quả của phép tính $815\,640 + 138\,785$ là:}{$954\,425$}{$953\,425$}{$944\,425$}{$954\,325$}
\cauTN{2}{Phép tính nào sau đây đúng?}{$a + b = b + a$}{$a - b = b - a$}{$a : b = b : a$}{$a + b = a - b$}
\cauTN{3}{Giá trị của biểu thức $a + b + c$ với $a = 64\,000, b = 26\,000, c = 30\,000$ là:}{$110\,000$}{$120\,000$}{$130\,000$}{$100\,000$}
\cauTN{4}{Số thích hợp điền vào chỗ chấm $9\,\text{tấn}\,150\,\text{kg} = \dots\,\text{kg}$ là:}{$915$}{$9\,150$}{$91\,500$}{$90\,150$}
\cauTN{5}{Số thích hợp điền vào chỗ chấm $7\,\text{m}^2\,56\,\text{dm}^2 = \dots\,\text{dm}^2$ là:}{$756$}{$7\,056$}{$7\,560$}{$78$}
\cauTN{6}{Khoảng thời gian $4\,\text{phút}\,50\,\text{giây}$ bằng bao nhiêu giây?}{$270\,\text{giây}$}{$290\,\text{giây}$}{$300\,\text{giây}$}{$310\,\text{giây}$}
\cauTN{7}{Tìm số hạng chưa biết: $x + 63\,280 = 98\,450$:}{$x = 35\,170$}{$x = 35\,270$}{$x = 161\,730$}{$x = 34\,170$}
\cauTN{8}{Điền dấu thích hợp: $75\,000 + 16\,000 \;\dots\; 60\,000 + 31\,000$}{$>$}{$<$}{$=$}{Không so sánh được}
\cauTN{9}{Số lớn nhất có $6$ chữ số là:}{$999\,999$}{$1\,000\,000$}{$987\,654$}{$900\,000$}
\cauTN{10}{Làm tròn số $964\,800$ đến hàng trăm nghìn ta được số:}{$900\,000$}{$960\,000$}{$1\,000\,000$}{$950\,000$}
\cauTNhaiDong{11}{Chu vi một khu vườn hình chữ nhật là $360\,\text{m}$. Nửa chu vi khu vườn là:}{$90\,\text{m}$}{$180\,\text{m}$}{$270\,\text{m}$}{$720\,\text{m}$}
\cauTNhaiDong{12}{Một thư viện trường học đợt 1 nhập về $1\,800$ cuốn sách, đợt 2 nhập về $2\,200$ cuốn sách. Tổng số cuốn sách thư viện nhập cả hai đợt là:}{$3\,800$}{$4\,000$}{$4\,200$}{$4\,500$}

\vspace{0.4mm}
\noindent\textbf{\large PHẦN II. TỰ LUẬN (7,0 điểm)}

\noindent\textbf{Bài 1 (2,0 điểm). Đặt tính rồi tính:}
\vspace{0.4mm}

\noindent
\begin{tabularx}{\linewidth}{|X|X|X|X|}
\hline
\oDatTinh[4]{a)}{$624\,835 + 197\,486$} & 
\oDatTinh[4]{b)}{$485\,392 + 376\,849$} & 
\oDatTinh[4]{c)}{$716\,384 + 89\,657$} & 
\oDatTinh[4]{d)}{$456\,273 + 543\,727$} \tabularnewline
\hline
\end{tabularx}

\newpage
% -------------------- TRANG 2: MẶT SAU --------------------
\noindent
\begin{tabularx}{\linewidth}{|X|c|}
\hline
\textbf{PHẦN TỰ LUẬN (TIẾP THEO) -- PHIẾU BÀI TẬP TUẦN 11} & \textbf{MÃ ĐỀ: 407} \\
\hline
\end{tabularx}

\vspace{0.8mm}

% BÀI 2: TÌM X (4 Ý)
\khungBaiTap{%
  \textbf{Bài 2 (1,5 điểm). Tìm $x$:}\par\vspace{0.6mm}
  \noindent
  \begin{minipage}[t]{0.485\linewidth}
    \textbf{a)} $x - 47\,690 = 85\,430$\par
    \dongke{3}
    \vspace{0.4mm}
    \textbf{c)} $98\,500 - x = 54\,390$\par
    \dongke{3}
  \end{minipage}\hfill
  \begin{minipage}[t]{0.485\linewidth}
    \textbf{b)} $x + 64\,280 = 97\,350$\par
    \dongke{3}
    \vspace{0.4mm}
    \textbf{d)} $x : 9 = 16\,428$\par
    \dongke{3}
  \end{minipage}
}

\vspace{0.8mm}

% BÀI 3: ĐỔI ĐƠN VỊ (4 Ý - 1 CỘT TỪ TRÊN XUỐNG)
\khungBaiTap{%
  \textbf{Bài 3 (1,0 điểm). Viết số thích hợp vào chỗ chấm:}\par\vspace{0.8mm}
  \noindent
  a) $8\,\text{tấn}\,450\,\text{kg} = \dotfill\,\text{kg}$\par\vspace{1.4mm}
  \noindent
  b) $6\,\text{m}^2\,88\,\text{dm}^2 = \dotfill\,\text{dm}^2$\par\vspace{1.4mm}
  \noindent
  c) $3\,\text{thế kỉ}\,65\,\text{năm} = \dotfill\,\text{năm}$\par\vspace{1.4mm}
  \noindent
  d) $5\,\text{phút}\,35\,\text{giây} = \dotfill\,\text{giây}$\par
}

\vspace{0.8mm}

% BÀI 4: BÀI TOÁN GIẢI (11 DÒNG KẺ)
\khungBaiTap{%
  \textbf{Bài 4 (1,5 điểm).} Một cửa hàng nhập về 18{,}5 kg gạo nếp. Mỗi ki-lô-gam gạo nếp giá 26 000 đồng. Cửa hàng bán hết số gạo với giá 31 000 đồng mỗi ki-lô-gam. Hỏi sau khi bán hết số gạo đó, cửa hàng lãi bao nhiêu tiền?
  \par\vspace{0.3mm}%
  \noindent\textit{Bài giải:}
  \dongke{11}
}

\vspace{0.8mm}

% BÀI 5: TOÁN NÂNG CAO (8 DÒNG KẺ)
\khungBaiTap{%
  \textbf{Bài 5 (1,0 điểm).} Tính bằng cách thuận tiện nhất:\par\vspace{0.6mm}\centerline{$K = 0{,}35 \times 8{,}7 + 0{,}35 \times 1{,}3$}
  \par\vspace{0.3mm}%
  \noindent\textit{Bài làm:}
  \dongke{8}
}

\newpage
\end{document}
